\medskip

\textbf{Affirmation 1 : Vraie}

$(d)$ est dirigée par le vecteur $\vect{u}\begin{pmatrix}-2\\0\\-6\end{pmatrix}$ et l'axe des ordonnées par $\vect{\jmath}$ de coordonnées $\begin{pmatrix} 0\\1\\0\end{pmatrix}$.

Les deux droites ont des vecteurs directeurs non colinéaires (ils sont même orthogonaux) donc les droites ne sont ni parallèles, ni confondues.

L'axe des ordonnées passe par l'origine du repère, de coordonnées $(0~;~0~;~0)$ et est dirigé par $\vect{\jmath}$ donc une représentation paramétrique de l'axe des ordonnées est :
\[\left\{\begin{array}{l c l}
x&=&0 \\
y&=&0 + s \\
z&=&0
\end{array} \quad, \text { où } s \in \R\right.\]

$(d)$ et l'axe des ordonnées seront donc coplanaires si et seulement si ces deux droites sont sécantes, c'est-à-dire si et seulement si le système suivant a une solution :

$\begin{aligned}
\left\{\begin{array}{lcl}
3-2t&=&0 \\
-1&=&s\\
2-6t&=&0
\end{array}\right.
	&\iff
\left\{\begin{array}{lcl}
2t&=&3 \\
6t&=&2
\end{array}\right.
	&\iff
\left\{\begin{array}{lcl}
t&=&1,5 \\
s&=&-1\\
t&=&\dfrac{1}{3}
\end{array}\right.
\end{aligned}$

Le système n'admet pas de solution, donc les droites n'ont pas de point commun.

Les droites ne sont ni parallèles, ni confondues, ni sécantes, par élimination, elles sont donc non coplanaires.

\bigskip

\textbf{Affirmation 2 : Vraie}

Soit $\mathcal{Q}$  plan passant par A et orthogonal à la droite $(d)$.

$(d)$ est dirigée par le vecteur $\vect{u}\begin{pmatrix}-2\\0\\-6\end{pmatrix}$

$\vect{u}\begin{pmatrix}-2\\0\\-6\end{pmatrix}$ est donc un vecteur normal au plan $\mathcal{Q}$  et donc aussi $\vect{n} \begin{pmatrix} 1\\0\\3 \end{pmatrix}$ car $\vect{u}$ et $\vect{n}$ sont colinéaires et non nuls.

Le plan $\mathcal{Q}$ admet donc une équation de la forme $x+ 0y + 3z + d = 0$ avec $d \in \R$. 

De plus, le point A appartient au plan $\mathcal{Q}$ donc $x_{\mathrm{A}}+3z_{\mathrm{A}}+d=0$ soit $3 +3\times (-2) +d=0$.

On a donc $d = -3 + 6 = 3$.

Une équation cartésienne du plan passant par A et orthogonal à la droite $(d)$  est donc : $x+3 z+3=0$.

\bigskip

\textbf{Affirmation 3 : Fausse}

Afin de déterminer l'angle $\widehat{\mathrm{BAC}}$ nous allons déterminer le produit scalaire $\vect {\mathrm{AB}} \cdot \vect {\mathrm{AC}}$

Déterminons tout d'abord les coordonnées du point C :

C est un point de $(d)$, il existe donc un réel $t$ tel que  : $\left\{\begin{array}{l}
x_{\mathrm{C}}=3-2t \\
y_{\mathrm{C}}=-1 \\
z_{\mathrm{C}}=2-6t
\end{array} \right.$

Or C est le point d'abscisse 2, on  a donc $x_{\mathrm{C}}=2$ d'où  $\begin{aligned}[t] 3-2t=2 & \iff 1= 2t \\ & \iff t=\dfrac12\end{aligned}$.

On a donc $z_{\mathrm{C}}=2- 6\times \dfrac{1}{2}= 2-3=-1$

Les coordonnées du point C sont donc $(2~;~-1~;~-1)$.

On peut donc calculer les coordonnées des vecteurs $\vect{\mathrm{AB}}$ et $\vect{\mathrm{AC}}$ : 

$\vect{\mathrm{AB}}\begin{pmatrix}x_\mathrm{B} -x_\mathrm{A} \\ y_\mathrm{B}-y_\mathrm{A}\\z_\mathrm{B}-z_\mathrm{A}\end{pmatrix}$ c'est-à-dire $\vect{\mathrm{AB}}\begin{pmatrix}2 \\ -1\\1\end{pmatrix}$ 
\quad de même : $\vect{\mathrm{AC}}\begin{pmatrix}-1 \\ 2\\1\end{pmatrix}$

Nous sommes dans un repère orthonormé, nous avons donc :

$\vect {\mathrm{AB}} \cdot \vect {\mathrm{AC}}= 2 \times (-1) + (-1) \times 2 + 1 \times 1 = -3$ 

Calculons les normes des  vecteurs $\vect{\mathrm{AB}}$ et $\vect{\mathrm{AC}}$ :

$\left\|\vect {\mathrm{AB}}\right\|^2=2^2+(-1)^2+1^2=6$, donc, comme une norme est positive $\left\|\vect {\mathrm{AB}}\right\|=\sqrt{6}$.

$\left\|\vect {\mathrm{AC}}\right\|^2=(-1)^2+2^2+1^2=6$, donc, de même :
$\left\|\vect {\mathrm{AC}}\right\|=\sqrt{6}$.

Or $\vect {\mathrm{AB}} \cdot \vect {\mathrm{AC}}=\mathrm{AB} \times \mathrm{AC} \times \cos\left(\widehat{\mathrm{BAC}}\right)$

on en déduit donc que :\quad $-1=\sqrt{6} \times  \sqrt{6} \times \cos \left(\widehat{\mathrm{BAC}}\right)$,

soit \quad$\cos \left(\widehat{\mathrm{BAC}}\right)=\dfrac{-3}{\sqrt{6}\sqrt{6}}=-\dfrac{3}{6}=-\dfrac{1}{2}$

et donc une mesure, exprimée en radian, de l'angle géométrique
$\widehat{\mathrm{BAC}}$ est $\dfrac{2\pi}{3}$.

\bigskip

\textbf{Affirmation 4 : Vraie}

Déterminons les coordonnées du point H.

H est le projeté orthogonal du point B sur la plan $\mathcal{P}$, c'est donc l'intersection du plan $\mathcal{P}$ et de la droite $\Delta$ orthogonal au plan $\mathcal{P}$ passant par le point B.

Le vecteur $\vect{u'}\begin{pmatrix}1\\0\\3\end{pmatrix}$ est un vecteur normal au plan $\mathcal{P}$, c'est donc un vecteur directeur de la droite $\Delta$.

De plus le point B$(5~;~-4~;~-1)$ appartient à la droite $\Delta$, une représentation paramétrique de la droite $\Delta$ est donc :
\[\left\{\begin{array}{lcr}
x&=&\phantom{-}5+t \\
y&=&-4 + 0t \\
z&=&-1+3t
\end{array} \quad, \text { où } t \in \R\right.\]

L'intersection entre la droite $\Delta$ et le plan $\mathcal{P}$ donne l'équation en la variable $t$ :

$\begin{aligned} (5+t)+3(-1+3t)-7=0 &\iff 5+t-3+9t-7=0\\
	& \iff 10t=5\\
	& \iff t=\dfrac{1}{2}
\end{aligned}$

Le point H d'intersection entre la droite $\Delta$ et $\mathcal{P}$ est le point de paramètre  $ t=\dfrac{1}{2}$ dans la représentation paramétrique de la droite $\Delta$.

Il a comme coordonnées :\quad$\mathrm{H}\left(5 + \dfrac{1}{2}~;~-4~;~-1 + 3 \times\dfrac{1}{2}\right)$ soit $\mathrm{H}\left(\dfrac{11}{2}~;~-4~;~\dfrac{1}{2}\right)$.

On a donc : $\vect{\mathrm{BH}}\begin{pmatrix}\frac{1}{2}\\0\\\frac{3}{2}\end{pmatrix}$

d'où : $\left\|\vect {\mathrm{BH}}\right\|^2=\left(\dfrac{1}{2}\right)^2+0^2+\left(\dfrac{3}{2}\right)^2=\dfrac{1}{4}+\dfrac{9}{4}=\dfrac{10}{4}$, donc $\left\|\vect {\mathrm{BH}}\right\|=\sqrt{\dfrac{10}{4}}=\dfrac{\sqrt{10}}{2}$.

La distance BH est égale à $\dfrac{\sqrt{10}}{2}$.
